% ============================================================
% PHẦN 3 — LỰA CHỌN MÔ HÌNH & KIẾN TRÚC
% ============================================================

\section{Lựa chọn mô hình và kiến trúc}
\label{sec:models}

Nhóm triển khai và so sánh \textbf{5 mô hình segmentation} thuộc hai nhóm kiến trúc: CNN-based và Vision Transformer-based (ViT-based). Mỗi mô hình được huấn luyện độc lập trên \emph{2 tập riêng biệt} (Rice và Coffee) trên Kaggle Notebook (GPU T4/P100).

\subsection{Tổng quan các mô hình}

\begin{table}[H]
    \caption{Tổng quan 5 mô hình segmentation được lựa chọn}
    \centering
    \begin{tabularx}{\textwidth}{llXl}
        \toprule
        \textbf{Mô hình} & \textbf{Loại} & \textbf{Đặc điểm nổi bật} & \textbf{Người phụ trách} \\
        \midrule
        Mask R-CNN       & CNN-based     & Baseline hai nhánh (detection + segmentation) & Lê Xuân Trí \\
        YOLO26-seg       & CNN-based     & Real-time, kiến trúc một nhánh & Nguyễn Hồ Anh Tuấn \\
        RF-DETR          & ViT-based     & End-to-end detection transformer (SOTA) & Tống Thanh Phúc \\
        MobileSAM        & ViT-based     & SAM rút gọn, segment anything & Đàm Tiến Đạt \\
        Mask2Former      & ViT-based     & Universal segmentation SOTA & Dương Tuấn Anh \\
        \bottomrule
    \end{tabularx}
    \label{tab:model_overview}
\end{table}

% ============================================================
% Hộp nhắc nhở yêu cầu — hiển thị một lần, áp dụng cho tất cả
% ============================================================
\begin{center}
\fbox{%
    \parbox{0.92\textwidth}{%
        \textbf{Yêu cầu bắt buộc — Kiến trúc chi tiết (áp dụng cho cả 5 mô hình):}\\[0.3em]
        Mỗi thành viên phụ trách mô hình phải hoàn thiện đầy đủ ba nội dung sau trong
        mục \emph{Kiến trúc chi tiết} của mô hình mình:
        \begin{enumerate}
            \item \textbf{Vẽ sơ đồ kiến trúc} --- Sơ đồ khối thể hiện luồng dữ liệu
                  qua từng thành phần; phải có chú thích hình bên dưới.
            \item \textbf{Mô tả số lượng tham số} --- Nêu tổng số tham số của mô hình
                  (trainable \& frozen) và số tham số từng thành phần chính.
            \item \textbf{Hàm kích hoạt được sử dụng} --- Liệt kê hàm kích hoạt tại
                  từng vị trí trong mạng (backbone, neck, head, decoder...).
        \end{enumerate}
        \textit{Thiếu bất kỳ điểm nào trong ba điểm trên là \textbf{chưa đạt yêu cầu}.}
    }%
}
\end{center}

% ----------------------------------------
\subsection{Mô hình 1: Mask R-CNN --- \textit{Lê Xuân Trí}}
\label{subsec:maskrcnn}

\subsubsection{Lý do lựa chọn}
Mask R-CNN \cite{he2017maskrcnn} là mô hình instance segmentation CNN nền tảng, được sử dụng làm \textbf{baseline} để so sánh với các mô hình ViT-based. Kiến trúc hai nhánh tách biệt giúp dễ debug và hiểu rõ từng thành phần.

\subsubsection{Kiến trúc chi tiết}

\begin{enumerate}
    \item \textbf{Sơ đồ kiến trúc} \textit{[Lê Xuân Trí điền]}\\
          \textit{[TODO: Chèn hình sơ đồ kiến trúc Mask R-CNN với chú thích bên dưới.
          Phải thể hiện rõ: Image → Backbone (ResNet-FPN) → RPN → RoI Align → Cls/Box/Mask heads.]}
    % \begin{figure}[H]
    %     \centering
    %     \includegraphics[width=0.85\textwidth]{img/maskrcnn_arch.png}
    %     \caption{Kiến trúc Mask R-CNN với backbone ResNet-50 FPN}
    %     \label{fig:maskrcnn_arch}
    % \end{figure}

    \item \textbf{Số lượng tham số} \textit{[Lê Xuân Trí điền]}\\
          \textit{[TODO: Điền tổng số tham số trainable/frozen, ví dụ:
          Backbone ResNet-50 FPN $\approx$~?M params; RPN $\approx$~?M; Mask head $\approx$~?M;
          Tổng: $\approx$~44M params (ước tính theo paper gốc --- cần xác nhận lại với
          \texttt{sum(p.numel() for p in model.parameters())}).]}

    \item \textbf{Hàm kích hoạt} \textit{[Lê Xuân Trí điền]}\\
          \textit{[TODO: Liệt kê hàm kích hoạt tại từng thành phần.
          Gợi ý: ReLU (ResNet backbone, FPN, RPN convs); Sigmoid (mask output head).
          Cần xác nhận lại theo code thực tế.]}
\end{enumerate}

% ----------------------------------------
\subsection{Mô hình 2: YOLO26-seg --- \textit{Nguyễn Hồ Anh Tuấn}}
\label{subsec:yolo}

\subsubsection{Lý do lựa chọn}
YOLO26-seg là phiên bản segmentation của YOLO --- kiến trúc CNN one-stage nổi tiếng về \textbf{tốc độ inference cao}. Nhóm chọn mô hình này để đánh giá trade-off giữa tốc độ và độ chính xác, đặc biệt quan trọng khi deploy trên web với yêu cầu latency $< 500$ms/ảnh.

\subsubsection{Kiến trúc chi tiết}

\begin{enumerate}
    \item \textbf{Sơ đồ kiến trúc} \textit{[Nguyễn Hồ Anh Tuấn điền]}\\
          \textit{[TODO: Chèn hình sơ đồ kiến trúc YOLO26-seg với chú thích bên dưới.
          Phải thể hiện rõ: Image → CSPDarknet backbone → PANet neck → Decoupled head (detect + segment).]}

    \item \textbf{Số lượng tham số} \textit{[Nguyễn Hồ Anh Tuấn điền]}\\
          \textit{[TODO: Điền tổng số tham số theo variant được chọn (nano/small/medium/large).
          Ghi rõ số params của backbone, neck, và detection/mask head riêng biệt.]}

    \item \textbf{Hàm kích hoạt} \textit{[Nguyễn Hồ Anh Tuấn điền]}\\
          \textit{[TODO: Liệt kê hàm kích hoạt tại từng thành phần.
          Gợi ý: SiLU/Swish (backbone, neck); Sigmoid (objectness, class scores, mask coefficients).
          Cần xác nhận theo code thực tế.]}
\end{enumerate}

% ----------------------------------------
\subsection{Mô hình 3: RF-DETR --- \textit{Tống Thanh Phúc}}
\label{subsec:rfdetr}

\subsubsection{Lý do lựa chọn}
RF-DETR \cite{rfdetr2024} là Detection Transformer của Roboflow --- mô hình \textbf{end-to-end} không cần NMS, thiết kế cho fine-tuning nhanh trên dataset nhỏ/trung bình. Là đại diện ViT-based SOTA trong nhóm.

\subsubsection{Kiến trúc chi tiết}

\begin{enumerate}
    \item \textbf{Sơ đồ kiến trúc} \textit{[Tống Thanh Phúc điền]}\\
          \textit{[TODO: Chèn hình sơ đồ kiến trúc RF-DETR với chú thích bên dưới.
          Phải thể hiện rõ: Image → DINOv2 backbone → Deformable Attention Encoder → Transformer Decoder → Bipartite matching output.]}

    \item \textbf{Số lượng tham số} \textit{[Tống Thanh Phúc điền]}\\
          \textit{[TODO: Điền tổng số tham số của RF-DETR.
          Ghi rõ params của DINOv2 encoder (frozen vs fine-tuned), decoder, và detection head.]}

    \item \textbf{Hàm kích hoạt} \textit{[Tống Thanh Phúc điền]}\\
          \textit{[TODO: Liệt kê hàm kích hoạt tại từng thành phần.
          Gợi ý: GELU (transformer FFN layers); Sigmoid (object classification).
          Cần xác nhận theo code thực tế.]}
\end{enumerate}

% ----------------------------------------
\subsection{Mô hình 4: MobileSAM --- \textit{Đàm Tiến Đạt}}
\label{subsec:mobilesam}

\subsubsection{Lý do lựa chọn}
MobileSAM \cite{mobile_sam} là phiên bản rút gọn của SAM (Segment Anything Model), thay thế ViT-H encoder bằng \textbf{TinyViT} --- nhỏ hơn 60× và nhanh hơn 40×, phù hợp deploy. Nhóm fine-tune để áp dụng vào domain bệnh lá cây.

\subsubsection{Kiến trúc chi tiết}

\begin{enumerate}
    \item \textbf{Sơ đồ kiến trúc} \textit{[Đàm Tiến Đạt điền]}\\
          \textit{[TODO: Chèn hình sơ đồ kiến trúc MobileSAM với chú thích bên dưới.
          Phải thể hiện rõ: Image → TinyViT Image Encoder → Prompt Encoder → Mask Decoder → Binary mask output.]}

    \item \textbf{Số lượng tham số} \textit{[Đàm Tiến Đạt điền]}\\
          \textit{[TODO: Điền tổng số tham số.
          Gợi ý tham khảo: TinyViT-5M encoder $\approx$~5.7M; Prompt encoder $\approx$~6K; Mask decoder $\approx$~3.87M; Tổng $\approx$~9.66M.
          Cần xác nhận lại theo code thực tế.]}

    \item \textbf{Hàm kích hoạt} \textit{[Đàm Tiến Đạt điền]}\\
          \textit{[TODO: Liệt kê hàm kích hoạt tại từng thành phần.
          Gợi ý: GELU (TinyViT transformer blocks); Sigmoid (mask output); LayerNorm (normalization).
          Cần xác nhận theo code thực tế.]}
\end{enumerate}

% ----------------------------------------
\subsection{Mô hình 5: Mask2Former}
\label{subsec:mask2former}

\subsubsection{Lý do lựa chọn}
Mask2Former \cite{cheng2022mask2former} thống nhất ba bài toán segmentation (instance, semantic, panoptic) trong một kiến trúc duy nhất dựa trên cơ chế \emph{mask classification}, thay vì dự đoán nhãn cho từng pixel như các mạng FCN truyền thống. Mô hình đạt kết quả SOTA trên COCO tại thời điểm công bố và là đại diện Transformer-based đối lập với baseline CNN (Mask R-CNN) trong nghiên cứu của nhóm. Việc fine-tune một kiến trúc Transformer SOTA trên bộ dữ liệu lá lúa và cà phê Việt Nam cũng chính là đóng góp sáng tạo S2 của nhóm: nó cho phép đánh giá khả năng chuyển giao (transferability) của các trọng số pre-trained trên COCO sang một domain hẹp, ít dữ liệu và đặc thù.

\subsubsection{Kiến trúc chi tiết}

\textbf{a) Sơ đồ kiến trúc}

Mask2Former gồm ba thành phần nối tiếp: một backbone trích đặc trưng đa tỷ lệ, một \emph{pixel decoder} nâng độ phân giải đặc trưng, và một \emph{Transformer decoder} sinh dự đoán theo từng truy vấn (query). Điểm cốt lõi tạo nên hiệu năng của mô hình là \emph{masked attention}: ở mỗi lớp decoder, mỗi query chỉ chú ý (attend) vào vùng foreground do chính nó dự đoán ở lớp trước, giúp hội tụ nhanh hơn và định vị đối tượng tốt hơn so với cross-attention toàn cục. Luồng dữ liệu được tóm tắt trong Hình~\ref{fig:m2f_arch}.

\begin{figure}[H]
    \centering
    \begin{tikzpicture}[
        node distance=0.55cm and 0.75cm,
        box/.style={draw, rounded corners, align=center, minimum height=1.05cm,
                    text width=2.5cm, font=\small, fill=blue!6},
        io/.style={draw, align=center, minimum height=1.05cm, text width=1.9cm,
                   font=\small, fill=gray!12},
        outbox/.style={draw, rounded corners, align=center, minimum height=1.05cm,
                    text width=2.2cm, font=\small, fill=green!8},
        arr/.style={-{Latex[length=2mm]}, thick}
    ]
        \node[io] (img) {Ảnh đầu vào\\$H\times W\times 3$};
        \node[box, right=of img] (bb) {Backbone\\Swin-T/S\\(đặc trưng đa tỷ lệ)};
        \node[box, right=of bb] (pd) {Pixel Decoder\\MSDeformAttn};
        \node[box, right=of pd] (dec) {Transformer Decoder\\$9$ lớp, masked attention\\$N{=}100$ query};
        \node[outbox, right=of dec] (cls) {Lớp bệnh\\(softmax)};
        \node[outbox, below=of cls] (msk) {Mask nhị phân\\(sigmoid)};

        \draw[arr] (img) -- (bb);
        \draw[arr] (bb) -- (pd);
        \draw[arr] (pd) -- (dec);
        \draw[arr] (dec) -- (cls);
        \draw[arr] (dec.south) |- (msk.west);
        \draw[arr] (pd.south) |- ([yshift=-0.18cm]msk.south west)
                   node[midway, below, font=\scriptsize] {per-pixel embedding};
    \end{tikzpicture}
    \caption{Sơ đồ kiến trúc Mask2Former. Backbone Swin sinh đặc trưng đa tỷ lệ; pixel decoder (multi-scale deformable attention) tạo embedding theo pixel và đặc trưng phân giải cao; Transformer decoder dùng $N$ query học được, qua $9$ lớp masked attention để sinh, với mỗi query, một nhãn lớp (softmax) và một vector mask embedding. Tích vô hướng giữa mask embedding và per-pixel embedding cho ra mask logits, kích hoạt bằng sigmoid.}
    \label{fig:m2f_arch}
\end{figure}

\textbf{b) Số lượng tham số}

Nhóm sử dụng hai backbone Swin Transformer ở hai cấu hình thử nghiệm. Tổng số tham số ước lượng theo checkpoint pre-trained của Hugging Face (\texttt{facebook/\allowbreak mask2former-\allowbreak swin-\{tiny,small\}-\allowbreak coco-instance}); con số chính xác có thể kiểm chứng bằng \texttt{sum(p.numel() for p in model.parameters())}. Backbone chiếm phần lớn dung lượng, trong khi pixel decoder và Transformer decoder tương đối nhẹ (Bảng~\ref{tab:m2f_params}).

\begin{table}[H]
    \caption{Phân bố tham số của Mask2Former theo hai backbone được sử dụng}
    \centering
    \begin{tabular}{lrr}
        \toprule
        \textbf{Thành phần} & \textbf{Swin-T (baseline)} & \textbf{Swin-S (tuned)} \\
        \midrule
        Backbone (Swin)          & $\approx$ 28M  & $\approx$ 50M  \\
        Pixel decoder (MSDeformAttn) & $\approx$ 10M & $\approx$ 10M \\
        Transformer decoder      & $\approx$ 9M   & $\approx$ 9M   \\
        \midrule
        \textbf{Tổng}            & $\approx$ 47M  & $\approx$ 69M  \\
        Dung lượng checkpoint (fp32) & 181{,}07 MB & 262{,}48 MB \\
        \bottomrule
    \end{tabular}
    \label{tab:m2f_params}
\end{table}

\textbf{c) Hàm kích hoạt}

Mô hình dùng \textbf{GELU} trong các khối MLP của Swin backbone và trong các lớp feed-forward của Transformer decoder, \textbf{ReLU} trong các lớp feed-forward của pixel decoder, kèm \textbf{LayerNorm} làm chuẩn hóa xuyên suốt các khối attention. Ở đầu ra, nhánh phân loại dùng \textbf{softmax} trên tập nhãn (cộng nhãn rỗng $\emptyset$), còn nhánh mask dùng \textbf{sigmoid} để đưa mask logits về xác suất foreground theo từng pixel.
